%! Quadratic Functions Revisited — Unit 2, Lesson 2.1 — Student Packet Cover Sheet
\documentclass[10pt]{article}
\usepackage{precalculus-article}
\usepackage{precalculus-boxes}
\usepackage{ltablex}
\keepXColumns

\begin{document}

% Full-bleed plum banner
\begin{tikzpicture}[remember picture, overlay]
  \fill[plum] (current page.north west)
    rectangle ([yshift=-0.9in]current page.north east);
\end{tikzpicture}
\vspace{-0.8in}

\begin{center}
  {\color{white}\textbf{\LARGE Precalculus}}\\[4pt]
  {\color{lilacmid}\large\textbf{Unit 2: Polynomial Functions}}\\[6pt]
  {\color{white}\textbf{\large Lesson 2.1 \quad Quadratic Functions Revisited}}
\end{center}

\vspace{0.22in}
\namedateperiod
\vspace{0.18in}

\begin{learningtargetbox}
\small
\begin{enumerate}[leftmargin=1.5em, itemsep=5pt]
  \item I can read the vertex, the axis of symmetry, and the zeros of a
        quadratic from its standard, vertex, and factored forms, and move
        from standard form to the vertex using $x = -\dfrac{b}{2a}$.
  \item I can compute the average rate of change of a linear or quadratic
        function over an interval, from a formula or a table, and recognize
        it as the slope of a secant line.
  \item I can build a table of first and second differences and use it to
        decide whether a function is linear or quadratic, to recover the
        leading coefficient, and to justify concavity from numerical
        evidence.
\end{enumerate}
\end{learningtargetbox}

\vspace{0.14in}

\begin{tocbox}
\small
\renewcommand{\arraystretch}{1.5}
\begin{tabularx}{\linewidth}{@{} c l X r @{}}
\rowcolor{lilac}
\textbf{\#} & \textbf{Component} & \textbf{Description} & \textbf{Score} \\
\midrule
1 & Warm-Up        & Spiral review of prerequisite skills               & \blank{1.2cm} \\
2 & Guided Notes   & Three forms, AROC, and difference tables          & \blank{1.2cm} \\
3 & Group Activity & Nova Cycles: pricing a helmet for revenue         & \blank{1.2cm} \\
4 & Exit Ticket    & Independent check on AROC and difference patterns & \blank{1.2cm} \\
5 & Homework       & Practice and extension                            & \blank{1.2cm} \\
\midrule
  & \multicolumn{2}{r}{\textbf{Total}} & \blank{1.2cm} \\
\end{tabularx}
\end{tocbox}

\vspace{0.14in}

\begin{remindbox}
\small
Lesson 2.0 introduced the polynomial family; today we return to its most
familiar member and look at it with Unit 1's tools. The parabola you graphed
in Algebra 2 has a second life: its average rates of change over equal steps
are not constant, but they change at a perfectly constant rate --- and that
single fact will identify a quadratic from nothing but a column of numbers.
\end{remindbox}

\end{document}
